Text guidance incorporates tracer descriptions and patient information (age, sex, weight and clinical scores) encoded by frozen BiomedCLIP encoders and injected into the RF model through cross-attention. The affine fine-tuning formulation optimizes lightweight adapters of BiomedCLIP-derived tracer guidance, rather than BiomedCLIP encoder weights, using a soft cross-tracer similarity penalty (Fig.~\ref{fig:1}B). In a separate adapter-level analysis, varying $\tau$ from 0.70 to 0.98 changed the mean cross-tracer embedding similarity from 0.7032 to 0.9807. At $\tau=0.82$, the similarity was 0.8214, compared with 0.9897 without the cross-tracer penalty (Table~\ref{tab:supp_tau_adapter}). These are representation-level results, not a comparison of PET generation quality across $\tau$. The reconstructed Full generation experiments used the original fixed tracer embeddings; their performance should not be attributed to affine adaptation.

A separate E1 fixed-weight sensitivity analysis replaced the clinical-score portion of the patient prompt while retaining age, sex and weight, MRI, tracer and model parameters (Table~\ref{tab:e1_score_removal}). Removing scores increased ROI MAE by approximately 0.000036 for FDG and 0.000079 for AV45, with statistical support under the endpoint-specific two-tracer Holm families but very small absolute changes. SSIM increased while PSNR decreased, and the average absolute prediction change was approximately 0.001019/0.001105 for FDG/AV45. These observations show conditional output sensitivity; they are not a retrained No-Scores ablation or proof of independent diagnostic benefit.
